\medskip

\emph{Cet exercice est un questionnaire à choix multiples.\\
Pour chacune des questions suivantes, une seule des quatre réponses proposées est exacte.\\ Une réponse exacte rapporte un point. Une réponse fausse, une réponse multiple ou l'absence de réponse à une question ne rapporte ni n'enlève de point.\\
Pour répondre, indiquer sur la copie le numéro de la question et la lettre de la réponse choisie.\\
Aucune justification n'est demandée.}

\medskip

\begin{enumerate}
\item On considère la fonction $f$ définie sur $\R$ par 
\[f(x) = x \e^{x^2 - 3}.\]

Une des primitives $F$ de la fonction $f$ sur $\R$ est définie par:
\begin{center}
\begin{tabularx}{\linewidth}{*{2}{X}}
\textbf{a.~}$F(x) = 2x\e^{x^2 - 3}$&\textbf{b.~}$F(x) = \left(2x^2 + 1\right)\e^{x^2 - 3}$\\
\textbf{c.~}$F(x) = \dfrac12 x\e^{x^2 - 3}$&\textbf{d.~}$F(x) = \dfrac12 \e^{x^2 - 3}$
\end{tabularx}
\end{center}

\item On considère la suite $\left(u_n\right)$ définie pour tout entier naturel $n$ par: \[u_n = \e^{2n+1}.\]

La suite $\left(u_n\right)$ est:

\begin{center}
\begin{tabularx}{\linewidth}{*{2}{X}}
\textbf{a.~~}arithmétique de raison 2 ; &\textbf{b.~~} géométrique de raison e;\\
\textbf{c.~~}géométrique de raison e$^2$ ; &\textbf{d.~~} convergente vers e. 
\end{tabularx}
\end{center}

Pour les questions \textbf{3.} et \textbf{4.}, on considère la suite $\left(u_n\right)$ définie sur $\N$ par:
\[u_0 = 15 \quad \text{et pour tout entier naturel }\: n \quad :\: u_{n+1} = 1,2u_n + 12.\]

\item La fonction Python suivante, dont la ligne 4 est incomplète, doit renvoyer la plus petite valeur de l'entier $n$ telle que $u_n > \np{10000}$.

\begin{center}
\begin{tabular}{|>{\ttfamily}l|}\hline
def seuil() :\\
\qquad n=0\\
\qquad u=15\\
\qquad while \ldots \ldots:\\
\qquad \quad n=n+1\\
\qquad \quad u=1,2$*$u+12\\
\qquad return(n)\\ \hline
\end{tabular}
\end{center}

À la ligne 4, on complète par :

\begin{center}
\begin{tabularx}{\linewidth}{*{4}{X}}
\textbf{a.~}u$\leqslant \np{10000}$ ; &\textbf{b.~}u $=\np{10000}$ &\textbf{c.~} u$> \np{10000}$;&\textbf{d.~}$n\leqslant \np{10000}$.\\
\end{tabularx}
\end{center}

\item On considère la suite $\left(v_n\right)$ définie sur $\N$ par : $v_n = u_n + 60$. 

La suite $\left(v_n\right)$ est:

\begin{center}
\begin{tabularx}{\linewidth}{*{2}{X}}
\textbf{a.~} une suite décroissante ; &\textbf{b.~} une suite géométrique de raison 1,2 ;\\
\textbf{c.~} une suite arithmétique de raison $60$ ;&\textbf{d.~} une suite ni géométrique ni arithmétique.
\end{tabularx}
\end{center}
\end{enumerate}

\bigskip
